\section{The Meantime}

Hope has the misfortune of a name worn smooth by weather forecasts, and the tradition's first service is to give it back its edges. The theological virtue of hope is not optimism, which is a temperament, nor wishing, which is idle. Its object is precise: a good that is \emph{future}, not yet possessed; \emph{arduous}, beyond the easy reach that needs no virtue; and \emph{possible}---but possible in a particular way, ``by means of the Divine assistance,'' so that in reaching for the good ``our hope attains God Himself, on Whose help it leans.''\endnote{\emph{Summa theologiae} II-II, q.~17, a.~1, corpus; checked 2026-07-25, English at New Advent, Latin at Corpus Thomisticum. The corpus analyzes hope's object as a future, difficult, possible good, and concludes that hope is a virtue precisely as leaning on the divine help (\latin{divino auxilio innitens} in the article's development). Hope, like charity, is infused, not acquired; the essay does not separately develop the psychology of the acquired hopes from which the theological virtue is distinguished.} \emph{Leans} is the load-bearing word. The posture of hope is not the sprinter's crouch, self-gathered for self-propulsion; the ceiling section retired that picture for good. It is the posture of a creature that has taken the measure of its arms, believed the promise anyway, and shifted its weight onto the one who made it---the vertigo's falling posture, turned by trust into something with a destination.

Sized against its object, the virtue becomes almost comically audacious. What may a creature hope for, from God? Aquinas answers by matching the hope to the helper rather than to the hoper: ``the good which we ought to hope for from God properly and chiefly is the infinite good, which is proportionate to the power of our divine helper\ldots{} Such a good is eternal life, which consists in the enjoyment of God Himself. For we should hope from Him for nothing less than Himself, since His goodness, whereby He imparts good things to His creature, is no less than His Essence.''\endnote{\emph{Summa theologiae} II-II, q.~17, a.~2, corpus (\latin{bonum quod proprie et principaliter a Deo sperare debemus est bonum infinitum\ldots{} vita aeterna, quae in fruitione ipsius Dei consistit, non enim minus aliquid ab eo sperandum est quam sit ipse}); checked 2026-07-25 as above. The proportion argued is to the power and goodness of the helper, not to the merits of the hoper---the exact reversal of every calculation the creature would make about itself.} \emph{Nothing less than Himself.} Every lesser hope aims at what the helper can give; this one aims at the giver, and its warrant is not the size of the creature's deserving---about which the entire series has been one long deflationary argument---but the size of the divine generosity, which is not smaller than the divine essence. Hope is the audit of section two, run in reverse: there, no created good was large enough to be the end; here, no created good is large enough to be the hope. The will's calibration and the virtue's object finally agree, and what they agree on is a person.

Between the two counterfeit postures, hope walks a ridge, and the counterfeits are old acquaintances of this series. On one side, presumption: the end grasped as if owed---spiritual bookkeeping's final form, in which the open account is imagined closed in the creature's favor and the vision awaited the way a creditor awaits payment. Part II diagnosed the desire underneath: to be done receiving, square at last, self-possessed in the one order where self-possession is a contradiction. On the other side, despair: the end refused as if the offer could not cover this case---which sounds like humility and is actually the old recoil in penitential dress, still insisting that the relation runs on the creature's adequacy, still refusing to let the gift be a gift. Both counterfeits make the same mistake in opposite directions: they take the creature's own performance as the measure of an end that was never measured by it. Hope differs from both not by moderation---it is far more extravagant than presumption, hoping for God from God---but by where it puts its weight. It leans; it does not bill, and it does not resign.\endnote{The mapping of presumption and despair onto Part II's ``spiritual bookkeeping'' and Part I's ``recoil'' is this essay's synthesis, labeled as such. The underlying doctrine---presumption and despair as the two vices opposed to hope, each corrupting trust in the divine help---is Aquinas's (ST II-II, qq.~20--21, treated summarily here and not separately quoted); nothing in this paragraph is a judgment about the culpability or interior state of any person, and pastoral treatment of despair lies outside this essay's scope.}

Paul's hymn to charity, read whole, gives the meantime its furniture and its horizon in six words apiece. ``We see now through a glass in a dark manner''---the mode of the interval, the mirror-knowledge this series has practiced from its first page. ``Charity never falleth away: whether prophecies shall be made void, or tongues shall cease, or knowledge shall be destroyed''---the demolition notice posted on everything proper to the interval as interval. ``And now there remain faith, hope, and charity, these three: but the greatest of these is charity.''\endnote{1 Corinthians 13:8, 12, 13, Douay--Rheims, checked 2026-07-25 as above. The traditional doctrinal reading summarized in this paragraph---faith and hope proper to the wayfaring state and ceasing in the vision they await, charity remaining because its act is the very substance of beatitude---is common theology grounded in the text's own contrasts; the essay does not enter the scholastic details of how faith and hope are ``evacuated'' in glory.} The tradition reads the ranking structurally, not sentimentally. Faith and hope are the virtues of the gap: faith believes what it cannot yet see, hope leans toward what it does not yet hold, and both are therefore self-retiring, scaffolding that comes down when the building stands. Charity is the building. The love infused now and the love of the blessed is one love, differing as the seed differs from the tree; it is the only part of the present life that is not provisional, the only act the creature performs in the dark that it will still be performing in the light. Which is why the greatest is charity: not because faith and hope are lesser duties, but because they are roads, and charity is the country.

So the meantime is not a waiting room, and the trilogy's earlier findings settle into their final places within it. The due return goes on---the knees still bend, the account stays open---but rendered now \emph{inside} a friendship, its acts no longer only justice's acknowledgment but love's correspondence: the letters written from the road. The vertigo goes on---the creature is as suspended at the end of this series as at its beginning---but the suspension has become the very thing hope leans into: being carried was never a predicament; it was transportation. What the interval adds is simply the direction of travel, and the tradition's last word about the destination is not an argument at all. It is a description---the most famous description of the end ever written from inside the dark manner of the mirror---and this series, like the book it comes from, can only finish by looking at it.
